%% sota_draft.tex -- Related Work / state-of-the-art draft (Phase 1.6).
%% Framed for IEEE TLT (learning analytics / EDM / assessment-tool register).
%% HARD CAP ~2 pages (CLAUDE.md S9.1): written before results exist; do not let it balloon.
%% \input this into main.tex (Phase 7). Citation keys resolve against article/references.bib.
%% ToE reframing notes are in % comments marked [ToE].
%% Rewritten 2026-10-16 (Pedro's review): the rubric subsection covers both domains, no criticism of prior
%% work, "overthinking" explained, the datasets/published-baselines subsection moved to Method (compact
%% datasets table), no method detail (parse rate) here.

\section{Related Work}
\label{sec:related}

Automated grading is a long-standing concern of the learning-analytics and educational-data-mining
communities~\cite{putnikovic_embeddings_2023,ferreira_mello_automatic_2025}. There, it is valued less
as a stand-alone scoring engine than as an assessment tool that must be reliable, transparent,
and fair before it can be used at scale for feedback or for grades that count~\cite{arz_von_straussenburg_enabling_2025,li_am_2023}.
Reliability in this sense has several faces (agreement with teachers, scoring error, and run-to-run
consistency), and the evidence reviewed below shows that a strong result on one does not guarantee the
others. Here we focus on what the most recent evidence implies for which configuration a grader
should use, the question this paper addresses.

\subsection{LLMs as Graders: Capable but Inconsistent}
Prompted LLMs now match or approach human raters on well-specified short-answer and essay
tasks~\cite{aminah_evaluating_2025,sessler_canai_2025,henkel_roars_2025}, yet their reliability is the
recurring caveat. Agreement is strong but depends on the model and the prompt, and consistency varies
sharply between systems: four LLMs grading the same algorithm answers with one rubric ranged from 0.29
to 0.83 in intraclass correlation (ICC), the standard index of agreement between repeated
measurements~\cite{moazzez_assessing_2024}. Repeated scoring also drifts over time~\cite{pack_large_2024},
and minimally-prompted grading can stay within one rubric point of teachers yet still show low exact
agreement with them~\cite{tate_can_2024}. Single scores are unstable enough that the variance across
repeated runs should be reported alongside the score~\cite{pack_large_2024}. Reliability matters beyond
measurement: grading errors can themselves harm learning, as false positives and negatives distort how
students revise their work~\cite{li_am_2023}. These observations motivate our repeated grading of every
answer, which measures consistency throughout.

\subsection{Reasoning in Grading: An Open Question}
The headline uncertainty concerns reasoning, the mode in which a model writes out its thinking before
it answers. In the general LLM-as-a-judge setting, explicit reasoning improves judgement: with a
thinking switch within one model family, Qwen3, ``thinking'' models gain roughly ten accuracy points at
modest extra cost~\cite{jayarao_explicit_2025}. The benefit depends on the task and the architecture: on
adjacent natural-language-generation evaluation (machine translation and summarisation), one reasoning
model improves with more reasoning while another underperforms its own non-reasoning
variant~\cite{larionov_deepseek_2025}. Studies of overthinking, where reasoning models write far longer
traces than a task needs, add that these traces carry verbose, redundant
computation~\cite{sui_stop_2025}. Grading a student answer is a different kind of task: the rubric bounds
what counts as correct, and the hard part is placing the score on the scale, which long deliberation may
not help. Whether reasoning helps or hurts here is therefore uncertain, and the three studies above
measure it in general judging or in adjacent generation tasks. Closest to grading, one study frames
short-answer grading as a move ``from memorization to reasoning'' and examines fine-tuning and
rubric-based prompting~\cite{cong_memorization_2026}. In essay scoring, two reasoning models agreed less
with expert raters, and less consistently, than a model without reasoning~\cite{yoshida_reasoning_2025},
but that comparison sets different models against each other. To our knowledge, no study has isolated
the effect of reasoning on grading itself, holding the model fixed and measuring quality, consistency,
and cost together. This paper does so with an on/off switch within each model (RQ1).

\subsection{Rubrics and Reference Answers}
What the grader is given to grade against matters. In short-answer grading the usual material is a
reference answer: the standard corpora pair each question with an instructor's or expert's answer, and
automatic graders have long scored a student answer by comparing it with that
answer~\cite{mohler_learning_2011,dzikovska_semeval_2013}. For LLM graders, the effect of a rubric has
been measured directly. In a comparison of ten LLMs, rubric-guided prompts improved agreement with human
essay scores~\cite{nugroho_evaluating_2025}, and a simplified rubric kept the essay-scoring accuracy of
three of four LLMs at the level of a detailed one while using fewer tokens~\cite{yoshida_detailed_2025}.
With GPT-3.5, a more explicit rubric made the grades of most of 18 engineering reports more consistent
across ten runs~\cite{garcia-huertes_leveraging_2024}, and closed GPT models scoring essays against ten
predefined criteria were more consistent and agreed better with teachers than open models, while rating
more leniently than the teachers~\cite{sessler_canai_2025}. In code grading, question-specific rubrics
outperform no-rubric prompting on object-oriented and data-structures code: Cohen's $\kappa$, a
chance-corrected agreement index, rises from 0.41--0.43 without a rubric to 0.60--0.65 with the best
question-specific rubric~\cite{pathak_rubric_2025}, and the best open-source code models grading
against instructor-specified rubrics reach a micro-accuracy of 85\%~\cite{havare_ai-based_2025}. Code
grading with LLMs is a younger and smaller area than short-answer grading: locally run open-source LLMs
already support teaching assistants in grading introductory programming~\cite{mazzone_exploring_2024},
and a systematic review of artificial intelligence in programming education lists automated assessment
and feedback among its core application areas~\cite{manorat_artificial_2025}.

How the model is told to apply the guidance has received less attention: one study compared a lenient
and a strict prompt for two LLMs grading geography worksheets~\cite{venkatraman_llmbased_2026}, and a
concurrent study found that a lenient prompt awarding partial credit lowered the grading error of every
model family on a mathematics exam~\cite{sapkota_math_2026}. Each of
these studies examines one kind of task. We measure the guidance and the instruction for applying it under
one protocol, on five models and four datasets of code and short answers.

\subsection{From Model Rankings to Configuration Guidance}
A growing body of work compares LLMs as graders: large side-by-side studies of programming
submissions~\cite{jukiewicz_comparison_2026}, assignments in a bioinformatics
course~\cite{policar_automated_2025}, and
multi-model frameworks emphasising transparency and reproducibility~\cite{anghel_graderassist_2025}.
Dedicated benchmarks for LLM short-answer scoring are also beginning to appear~\cite{lai_sasbench_2026}.
Where these studies compare which model grades best, we ask which configuration a grading task calls
for. Our contribution is a \emph{configuration guide}: an empirically grounded recommendation of how to
configure an LLM grader for computer-science exams, with the evidence and the cost behind each setting.

% [ToE] If the venue becomes IEEE Transactions on Education, re-foreground the pedagogical contribution:
% lead with instructor workload, assessment validity, fairness, and the computing-education context;
% expand the human-in-the-loop and false-positive-harm argument (cf. \cite{li_am_2023}); keep the
% model/benchmark detail terser. The study itself is unchanged.
